\begin{document}

\AddToShipoutPicture*{\BackgroundPic}

% --- CAPA ---
\clearpage
\thispagestyle{empty}

\begin{bfseries}
\begin{center}

\vspace{0.5cm} 
UNIVERSIDADE FEDERAL DO PARANÁ \\

\vspace{0.6cm}
SAMUEL KUTZ PARANHOS \\
\vspace{2.5cm}

\vspace*{1cm}
\textbf{\MakeUppercase{Scientific Machine Learning Methods for the \\ Nonlinear Boussinesq System of Equations}} \\

\begin{minipage}[t]{0.45\textwidth} \hspace{2cm}
\raggedright
\end{minipage}
\begin{minipage}[t]{0.5\textwidth}\hspace{2cm}
\raggedright
\vspace*{3cm}
\begin{flushleft}
\footnotesize  
\justify
\hspace{-0,35cm}
\textnormal{}
\end{flushleft}
\end{minipage}

\vspace*{3cm}
\textbf{\MakeUppercase{CURITIBA}} \\
\textbf{\MakeUppercase{\textnormal{2025}}}
\end{center}
\end{bfseries}

\pagebreak
% --- RESUMO ---
\chapter*{ABSTRACT}
Partial Differential Equations (PDEs) are the essential language of scientifc computing and many predictions, 
studies and models are based on solving some kind of differential equation. In this work, we study the efficacy of 
Scientific Machine Learning methods (SciML) for the Boussinesq System of PDEs, a shallow wave model based off the Euler Equations, 
from classical Hydrodynamics. This study focused on understanding the relationship between the PDE parameters, its solution and the 
SciML method performance, implementing Physics-Informed Neural Networks (PINNs) with and without data (semi-supervised) and 
Fourier Neural Operators (FNOs) (purely supervised), analyzing through comparisons to the reference solutions coming from 
pseudospectral classical discrete methods. We already understood that PINNs without data suffer from Spectral Bias, the tendency 
of neural networks of first learning the low-frequency components of the solution during the optimization process, thus making 
small details hard to learn. On the other hand, FNOs did not show Spectral Bias, even though it is known to exibit it too. (still more to do)

\vspace{1cm}
\noindent
\textbf{Keywords:} Scientific Machine Learning; Partial Differential Equations; Boussinesq System

\pagebreak

\tableofcontents
\pagebreak

% --- CORPO DO TEXTO ---

\chapter{INTRODUCTION}





Since the start of the scientific thinking, mathematics has been one of the key
languages in which we can understand the world, and, Partial Differential
Equations (PDEs) are one of the main dialects we use to model and to simplify
nature's complexity.

These equations do not always have simple closed-form solutions (that is, when
they have), showing the need of approximated numerical solutions coming from
numerical methods for PDEs in order to extract predictions, compare with data,
build simulations and so on. These classical methods, widely studied for all
types of equations, rely on discretizing the PDE's domain into grids and performing time steps.

With the ascension of Machine Learning (ML) methods and the whole ecosystem of
tools built around the field, a natural extesion is to build a bridge
connecting numerical solution of PDEs with ML methods, a field called Scientific
Machine Learning (SciML). These methods have been applied to many different types of PDEs \cite{SAKDASJD}, 
our focus is in its application to Hydrodyanamics.


\subsection{Literature Overview on Scientific Machine Learning}

Apresentado por Raisis...

Although SciML in a new field, the literature in this field is too vast, therefore, here, 
we focus on the review of the application of such methods restricted to classical Hydrodynamics equations. 

++lista de eq + methods



As presented in the literature overview, as far as we know, no study has investigated 
the efficiency of SciML methods on the Boussinesq System.

\subsection{The Boussinesq System of Equaitions}

quem é a equacao, de onde vem (mestrado luiz) e pq é relevante.

\subsection{The Research Goals and }

In order to fill this gap in the literature, our goal is to implement and evaluate the performance of two
famous models of SciML: the Physics-Informed Neural Networks (PINN) and the Fourier
Neural Operators (FNO) and Physics-Informed Neural Operators (PINO), which together represent a novel family of methods in
SciML, investigating the impact of data and the PDE model (so called
Physics-Informed) presence in training, revealing their advantages and
limitations for each method.

In order to achieve these goals, we implemented these methods in Python taking advantage of Pytorch, comparing with a reference solution obtained by pseudospectral numericla methods \cite

++ itemize com melhor organizacao das ideias e incluir PINO.

++ paragrafo de relevancia do trabalho

\subsection{Work Structure}

Apresentar capitulos, this work is written with the goal of sharing SciML 
for senior undergrad, we recommend the reader to have some knowlodge of scientifc computing and 

with this in mind, we organize this work as etc
 

\printbibliography

\end{document}